\documentclass[11pt]{article}
\usepackage[margin=0.55in]{geometry}
\usepackage{setspace}
\usepackage{booktabs}
\usepackage{longtable}
\usepackage{array}
\usepackage{amsmath, amssymb}
\usepackage{mathtools}
\usepackage{hyperref}
\usepackage{enumitem}
\usepackage{float}

\hypersetup{
  colorlinks=true,
  linkcolor=blue,
  citecolor=blue,
  urlcolor=blue
}

\pdfstringdefDisableCommands{%
  \def\beta{beta}%
  \def\alpha{alpha}%
  \def\delta{delta}%
  \def\gamma{gamma}%
  \def\tau{tau}%
  \def\sigma{sigma}%
  \def\varepsilon{epsilon}%
  \def\mathbb#1{#1}%
  \def\mathrm#1{#1}%
}

\newcommand{\E}{\mathbb{E}}
\newcommand{\Var}{\mathrm{Var}}
\newcommand{\Cov}{\mathrm{Cov}}
\newcommand{\R}{\mathbb{R}}

\begin{document}
\begin{center}
{\LARGE \textbf{Media Bias ISM}}\\
\vspace{0.15em}
{\normalsize Regressions Compendium --- All Specifications, Equations, and Results}\\
\end{center}
\vspace{0.2em}
\hrule
\vspace{0.5em}
\setstretch{1.05}
\small

\section*{Project Motivation}\label{sec:motivation}
The 2020 recession is a clean, universal, exogenous shock to every large-cap US firm. It removes the usual hard problem of finding a valid control group: every firm is treated simultaneously. That makes it a natural setting for testing pre/post structural changes in firm-level outcomes.

The project asks a chain of two questions:
\begin{enumerate}[leftmargin=1.7em]
\item \textbf{RQ1.} Did the \emph{tone} of media coverage of S\&P 500 firms shift after Jan 2020?
\item \textbf{RQ2.} Did firm-level \emph{returns} shift after Jan 2020 beyond what market co-movement explains?
\end{enumerate}

\section*{Data and Variable Construction}\label{sec:data}
\textbf{Corpus and sampling.} The raw scrape contains 4{,}706{,}589 firm--article links across 30 S\&P 500 firms, 2015--2025. We filter to articles whose titles explicitly contain the company name (not ticker, to avoid Apple-the-fruit / Visa-the-immigration problems), de-duplicate by normalized title within firm, and drop firms with more than three years of $\le 3{,}000$ articles. The retained set is 26 firms and 690{,}586 firm--article links.

\textbf{Daily stance series.} For each retained article $a$ of firm $i$ on date $t$, with NewsMTSC probabilities $(p_a^{+}, p_a^{-}, p_a^{0})$:
\[
s_a = p_a^{+} - p_a^{-},\qquad
\text{bias}_{it} = \frac{1}{|\mathcal{A}_{it}|}\sum_{a \in \mathcal{A}_{it}} s_a.
\]
$|\mathcal{A}_{it}|$ is the firm-day article count, also retained as a control $\text{vol}_{it}$. At $\tau=0$ the RQ1 panel has 63{,}200 firm-day observations.

\textbf{Daily returns and market controls.}
\[
\text{return}_{it}=\frac{P_{it}-P_{i,t-1}}{P_{i,t-1}},\qquad
\text{sp500\_return}_t = \frac{P^{\text{SP500}}_t - P^{\text{SP500}}_{t-1}}{P^{\text{SP500}}_{t-1}},\qquad
\text{vix}_t = \text{CBOE VIX close on day } t.
\]
The RQ2 panel has 62{,}667 firm-trading-day observations.

\textbf{Shock indicator.}
$
\text{Post}_t = \mathbb{1}[t \ge 2020\text{-}01\text{-}01],\quad
\text{pre: } 2015\text{-}2019,\quad
\text{post: } 2020\text{-}2025.
$

\begin{table}[H]
\centering
\caption*{Panel sizes across the three research questions}
\begin{tabular}{lrrl}
\toprule
Track & Observations & Firms & Date / frequency \\
\midrule
RQ1 (daily stance) & 63{,}200 firm-days & 26 & 2015-01-01 to 2025-12-31 daily \\
RQ2 (daily returns) & 62{,}667 firm-days & 26 & 2015-12-01 to 2025-11-24 daily \\
% RQ3 rows removed
\bottomrule
\end{tabular}
\end{table}

%================================
\section*{RQ1.1 --- Two-Way FE (Firm + Day) on Stance}\label{sec:rq1.1}
\textbf{Motivation.} Simplest pre/post test on stance. Add day fixed effects on top of firm fixed effects so that the Post coefficient is identified only from within-firm-within-day variation, net of every common shock that hit all firms on a given day.

\textbf{Specification.}
\[
\text{bias}_{it} = \alpha_i + \delta_t + \beta\,\text{Post}_t + \gamma\,\text{vol}_{it} + \varepsilon_{it}.
\]

\textbf{Sample.} $n = 63{,}200$ firm-days, 26 firms, 4{,}018 distinct trading days, $\tau = 0$.

\begin{table}[H]
\centering
\caption*{Table RQ1.1: Two-Way FE on stance (HC1 robust SE)}
\begin{tabular}{lrrrrrr}
\toprule
Variable & Coef. & SE (HC1) & $t$-stat & $p$-value & CI Low & CI High \\
\midrule
Post & 0.014130 & 0.013623 & 1.0372 & 0.2996 & $-$0.012571 & 0.040831 \\
Article Volume & $-$0.000970 & 0.000076 & $-$12.7724 & $<10^{-37}$ & $-$0.001119 & $-$0.000822 \\
\midrule
Within $R^2$ & 0.001749 & \multicolumn{5}{r}{} \\
\bottomrule
\end{tabular}
\end{table}

\textbf{Interpretation.} The post-2020 within-firm shift in stance is positive ($+0.014$) but \emph{not} significant ($p = 0.30$) after absorbing daily common shocks. Article volume is strongly negative and highly significant: one extra firm-day article is associated with a $0.001$ drop in stance, consistent with negative-bias-driven coverage spikes. Within $R^2$ is small (0.0017): most variance in $\text{bias}_{it}$ is firm-by-day idiosyncratic noise that two regressors cannot capture.

\subsection*{Pending: log-normalised volume}
The volume control is currently included in \emph{raw count} form. A log transform $\widetilde{\text{vol}}_{it} = \ln(1 + \text{vol}_{it})$ is the standard fix for right-skewed counts. It compresses the influence of high-coverage outliers (max 539 articles per firm-day) and converts the coefficient interpretation to ``percent change in stance per percent change in coverage.'' \emph{Open to-do --- not yet re-estimated.}

%================================
\section*{RQ1.2 --- Same Spec, Day-Clustered SE}\label{sec:rq1.2}
\textbf{Motivation.} HC1 assumes independent observations conditional on covariates. In a daily panel, two firms on the same day share macro shocks, so residuals can be correlated within day. Day-clustered SEs are robust to this dependence and are the preferred uncertainty estimate for daily panels.

\textbf{Specification.} Identical to RQ1.1, only the variance estimator changes:
\[
\widehat{\Var}(\hat{\theta}) = (\tilde{X}^{\top}\tilde{X})^{-1} \left(\sum_{g \in \text{days}} \tilde{X}_g^{\top}\hat{u}_g\hat{u}_g^{\top}\tilde{X}_g\right) (\tilde{X}^{\top}\tilde{X})^{-1}.
\]

\begin{table}[H]
\centering
\caption*{Table RQ1.2: Day-clustered SE (4{,}018 clusters)}
\begin{tabular}{lrrrrrr}
\toprule
Variable & Coef. & SE (clustered) & $t$-stat & $p$-value & CI Low & CI High \\
\midrule
Post & 0.014130 & 0.013432 & 1.0520 & 0.2928 & $-$0.012196 & 0.040456 \\
Article Volume & $-$0.000970 & 0.000076 & $-$12.8438 & $<10^{-37}$ & $-$0.001119 & $-$0.000822 \\
\bottomrule
\end{tabular}
\end{table}

\textbf{Interpretation.} Day-clustered SEs are virtually identical to HC1 SEs $(<1\%$ deviation$)$. Inference does not depend on the cluster assumption; the insignificance of Post is not an artifact of ignoring within-day correlation.

%================================
\section*{RQ1.3 --- Joint $F$-Test for Day Fixed Effects}\label{sec:rq1.3}
\textbf{Motivation.} The single Post dummy assumes the post-2020 shift is constant across all days. The day FE jointly span every individual daily shock. If they are not jointly significant, including them is overkill; if they are, the time variation matters and the static Post compresses dynamic signal.

\textbf{Specification.} $F$-test of $H_0: \delta_1 = \delta_2 = \dots = \delta_T = 0$ in the two-way FE model.

\begin{table}[H]
\centering
\caption*{Table RQ1.3: Joint significance of day FE}
\begin{tabular}{lr}
\toprule
Statistic & Value \\
\midrule
$F$-stat & 1.1348 \\
df1 (time periods $-$ 1) & 4{,}017 \\
df2 (residual df) & 59{,}181 \\
$p$-value & $1.12 \times 10^{-8}$ \\
$F$ critical (5\%) & 1.0383 \\
\bottomrule
\end{tabular}
\end{table}

\textbf{Interpretation.} Day FE are jointly significant at $p < 10^{-7}$. Daily common shocks matter and cannot be collapsed into a single Post indicator. This is the technical justification for the event-study specification in RQ1.4.

%================================
\section*{RQ1.4 --- Monthly Event Study on Stance}\label{sec:rq1.4}
\textbf{Motivation.} The static Post compresses everything post-Jan-2020 into one number. An event study replaces it with one coefficient per calendar month relative to Jan 2020, giving (i) a pre-trend test (months before should hover around zero) and (ii) the dynamic post-shock path.

\textbf{Specification.} Firm fixed effects + monthly event-time dummies, $k \in [-24, +24]$, reference $k = -1$ (Dec 2019):
\[
\text{bias}_{it} = \alpha_i + \sum_{k \neq -1} \beta_k \mathbb{1}[t \in \text{month}_k] + \gamma\,\text{vol}_{it} + \varepsilon_{it}.
\]
$n = 23{,}439$ firm-day observations within the $\pm 24$-month window around Jan 2020.

\begin{table}[H]
\centering
\caption*{Table RQ1.4: Selected event-study coefficients (HC1 SE)}
\begin{tabular}{lrrrrl}
\toprule
$k$ (month) & Coef. & SE & $t$-stat & $p$-value & Comment \\
\midrule
$-24$ & 0.04826 & 0.01791 & 2.694 & 0.007 & pre-trend outlier (Jan 2018) \\
$-12$ & 0.02169 & 0.01805 & 1.202 & 0.229 & pre-trend OK \\
$-6$  & 0.03176 & 0.01720 & 1.846 & 0.065 & pre-trend mild outlier \\
$-3$  & 0.03442 & 0.01704 & 2.020 & 0.043 & pre-trend mild outlier \\
$-1$  & 0       & ---     & ---   & ---   & reference (Dec 2019) \\
$\mathbf{0}$ (Jan 2020) & 0.03103 & 0.01835 & 1.691 & 0.091 & event month \\
$+3$  & 0.04215 & 0.01917 & 2.199 & 0.028 & early post-shock \\
$+6$  & 0.03220 & 0.01825 & 1.764 & 0.078 & mid post-shock \\
$+12$ & 0.06405 & 0.01880 & 3.408 & $<10^{-3}$ & 1-year post: significant \\
$+18$ & 0.05550 & 0.01878 & 2.955 & 0.003 & sustained \\
$+24$ & 0.05039 & 0.01871 & 2.694 & 0.007 & 2-year post: sustained \\
\bottomrule
\end{tabular}
\end{table}

\textbf{Interpretation.} Pre-period coefficients are mostly insignificant with a few mild outliers (Jan 2018, mid-2019) that do not form a systematic trend: parallel trends is reasonable. The event month $k = 0$ shows a marginal effect ($p = 0.09$). Effects emerge by $k = +3$ months (early COVID waves) and become \emph{persistently} positive by $k = +12$, staying at $+0.05$ to $+0.10$ through $k = +24$. The static Post (RQ1.1) is insignificant because it averages an absent month-0 effect with a strong post-year-1 effect; once dynamics are allowed, the post-2020 stance shift is clear and persistent.

%================================
\section*{RQ1.5 --- MLR / Gauss-Markov Diagnostics}\label{sec:rq1.5}
\textbf{Motivation.} Verify the assumptions under which OLS is BLUE; flag any failures that require switching estimators or SEs.

\begin{table}[H]
\centering
\caption*{Table RQ1.5: Diagnostic tests on the two-way FE residuals ($n = 63{,}200$)}
\begin{tabular}{lll}
\toprule
Assumption / Test & Result & Implication \\
\midrule
Residual mean & $-9.0 \times 10^{-19}$ & $\checkmark$ Orthogonality satisfied \\
Residual SD & 0.278 & $\checkmark$ Non-trivial variation \\
Condition number & 157.46 & $\checkmark$ Numerically stable \\
\midrule
\multicolumn{3}{l}{\textit{Heteroscedasticity (Breusch-Pagan)}} \\
LM statistic & 190.25 & $\times$ Heteroscedasticity present \\
$p$-value & $< 10^{-42}$ & HC1 robust SE used \\
\midrule
\multicolumn{3}{l}{\textit{Normality (Jarque-Bera)}} \\
JB statistic & 4{,}541.18 & $\times$ Non-normality \\
Skewness, Kurtosis & $-$0.198 / 4.252 & fat tails (typical for financial data) \\
\midrule
\multicolumn{3}{l}{\textit{Multicollinearity (VIF)}} \\
Post, Article Volume & 1.0004 each & $\checkmark$ no collinearity \\
Max pairwise correlation & 0.0189 & $\checkmark$ negligible \\
\midrule
\multicolumn{3}{l}{\textit{Serial correlation (Durbin-Watson)}} \\
DW statistic & 1.673 & moderate; partly absorbed by FE \\
\bottomrule
\end{tabular}
\end{table}

\textbf{Interpretation.} Two assumptions fail: residuals are heteroscedastic and non-normal. Both are handled by HC1 robust SEs and the large sample size ($n = 63{,}200$, well within CLT range). No collinearity. Residual mean is machine zero, confirming FE absorption. Day-clustered SEs (RQ1.2) virtually equal HC1 SEs, so within-day correlation is not large. Inference is sound.

%================================
\section*{RQ2.1 --- Baseline Firm FE on Returns}\label{sec:rq2.1}
\textbf{Motivation.} Mirror RQ1 on the returns side. If markets reacted to the same shock, we should see a within-firm pre/post shift in returns. Market and volatility controls strip out everything common to all firms.

\textbf{Specification.}
\[
\text{return}_{it} = \alpha_i + \beta\,\text{Post}_t + \gamma_1\,\text{sp500\_return}_t + \gamma_2\,\text{vix}_t + \varepsilon_{it}.
\]
$n = 62{,}667$ firm-trading-days. Pre-period firm means tightly clustered (range 0.00193, SD 0.00042), so the DiD parallel-trends assumption holds at the firm-mean level; no SCDiD fallback required.

\begin{table}[H]
\centering
\caption*{Table RQ2.1: Entity FE regression on returns (HC1 robust SE)}
\begin{tabular}{lrrrrrr}
\toprule
Variable & Coef. & SE & $t$-stat & $p$-value & CI Low & CI High \\
\midrule
Post & $-$0.000344 & 0.000172 & $-$2.0038 & 0.0451 & $-$0.000681 & $-$0.000008 \\
S\&P 500 Return & 1.127576 & 0.010982 & 102.677 & $\approx 0$ & 1.106052 & 1.149100 \\
VIX & $-$0.000001 & 0.000020 & $-$0.043 & 0.965 & $-$0.000039 & 0.000038 \\
\bottomrule
\end{tabular}
\end{table}

\textbf{Interpretation.} Post-2020 daily returns are about 3.4 basis points lower on average than pre-2020, marginally significant ($p = 0.045$) after stripping out market index movement and volatility. The S\&P 500 coefficient of about 1.13 is a textbook market beta. VIX is essentially zero once the index return is in (the two are highly correlated and the index absorbs the volatility-regime signal).

%================================
%================================
% RQ3 (sections RQ3.1 and RQ3.2) removed per user request.
%================================

%================================
\section*{Cross-Regression Summary}\label{sec:summary}
\begin{table}[H]
\centering
\caption*{All eight regressions at a glance}
\renewcommand{\arraystretch}{1.1}
\begin{tabular}{lp{0.18\linewidth}p{0.12\linewidth}p{0.30\linewidth}p{0.18\linewidth}}
\toprule
ID & Spec & Outcome & Headline result & Verdict \\
\midrule
RQ1.1 & Two-way FE + Post + vol & stance & $\beta(\text{Post}) = +0.014$, $p = 0.30$ & n.s.\ (static) \\
RQ1.2 & + day-clustered SE & stance & $\beta(\text{Post})$ same, $p = 0.29$ & robust to clustering \\
RQ1.3 & $F$-test on day FE & --- & $F = 1.13$, $p = 1.1 \times 10^{-8}$ & day FE matter \\
RQ1.4 & Monthly event study & stance & Persistent post-shock uplift: $\beta_{+12} = +0.064$, $p < 10^{-3}$ & sig.\ (dynamic) \\
RQ1.5 & MLR diagnostics & residuals & BP $p < 10^{-42}$; JB $p \approx 0$; VIF $\approx 1$ & HC1 SE justified \\
RQ2.1 & Firm FE + Post + market + VIX & return & $\beta(\text{Post}) = -0.000344$, $p = 0.045$ & marginally sig. \\
% RQ3 rows removed from summary table
\bottomrule
\end{tabular}
\end{table}

\textbf{Reading.} Static pre/post tests (RQ1.1, RQ2.1) are weak or marginal because they compress five post-shock years into a single dummy. Event-study specs (RQ1.4) tell the real story: stance shifts gradually and locks in higher by year 1; media reliably reacts to extreme price moves.

%================================
\section*{Identification and Robustness}\label{sec:robustness}
\textbf{Parallel trends.} Pre-period coefficients in the RQ1.4 event study are mostly small and insignificant. For RQ2.1, pre-period firm means are tightly clustered (SD $= 0.0004$), so the DiD assumption holds at the firm-mean level; no SCDiD fallback required.

\textbf{Cluster vs HC1.} Day-clustered SEs for RQ1.1 are within 1\% of HC1 SEs, so within-day cross-firm correlation is not driving inference.

\textbf{Heteroscedasticity and non-normality.} Both present (RQ1.5). Handled with HC1 robust covariance and large-sample CLT inference; reported $p$-values use the two-sided normal approximation $p = 2\Phi(-|t|)$.

\textbf{Day FE absorption.} The $F$-test (RQ1.3) confirms day FE are jointly significant; they are the default in RQ1. They are not currently used in RQ2.1; the same two-way FE upgrade is on the to-do list.

\textbf{Confidence threshold $\tau$.} Default $\tau = 0$. Sensitivity at $\tau \in \{0.6, 0.9\}$ is available via the RQ1 pipeline.

\textbf{Relevance filtering.} Before NewsMTSC stance scoring, articles pass through a DeBERTa-v3-base relevance classifier with operating point chosen to satisfy precision $\ge 0.90$ (ensemble current best: $P = 0.944$, $R = 0.810$ on a 187-article held-out test set).

%================================
% Open To-Dos removed per user request.
%================================
\section*{Variables Index}\label{sec:vars}
\begin{table}[H]
\centering
\caption*{All symbols used in this document}
\begin{tabular}{lll}
\toprule
Symbol & Meaning & First appears \\
\midrule
$i$ & Firm index ($i \in \{1, \ldots, 26\}$) & Data section \\
$t$ & Date index (daily) & Data section \\
$\alpha_i$ & Firm fixed effect & RQ1.1 \\
$\delta_t$ & Day (or month) fixed effect & RQ1.1 \\
$\beta$ & Main coefficient on Post indicator & RQ1.1 \\
$\beta_k$ & Event-study coefficient at relative time $k$ & RQ1.4 \\
$\gamma, \gamma_1, \gamma_2$ & Control-variable coefficients & RQ1.1 \\
$\varepsilon_{it}$ & Idiosyncratic error term & RQ1.1 \\
$\text{bias}_{it}$ & Daily firm-day stance score (NewsMTSC) & Data section \\
$s_a$ & Article-level stance: $p_a^{+} - p_a^{-}$ & Data section \\
$p_a^{+}, p_a^{-}, p_a^{0}$ & NewsMTSC probabilities (pos / neg / neu) & Data section \\
$\tau$ & NewsMTSC confidence threshold ($0$, $0.6$, $0.9$) & Glossary \\
$\text{vol}_{it}$ & Article count for firm $i$ on day $t$ & RQ1.1 \\
$\text{Post}_t$ & $\mathbb{1}[t \ge 2020\text{-}01\text{-}01]$ & Data section \\
$\text{return}_{it}$ & Daily simple return for firm $i$ & Data section \\
$\text{sp500\_return}_t$ & Daily S\&P 500 index return & Data section \\
$\text{vix}_t$ & CBOE VIX close on day $t$ & Data section \\
$n$ & Estimation sample size & all regressions \\
$K$ & Number of regressors & Glossary \\
$\widehat{SE}$ & Standard error of an estimated coefficient & Glossary \\
$\widehat{\Var}$ & Estimated variance-covariance matrix & Glossary \\
\bottomrule
\end{tabular}
\end{table}

\vspace{0.5em}
\hrule
\vspace{0.3em}
{\footnotesize \textit{Source files.}
RQ1 pipeline: \texttt{rq/rq-1-pipeline.py}.
RQ2 pipeline: \texttt{rq/rq-2-pipeline.py}.
MLR diagnostics: \texttt{rq/mlr\_assumptions\_check.py}.
STMA anomaly detection: \texttt{stma/anomaly\_det/}.
	extit{Result files.} \texttt{results/rq1/}, \texttt{results/rq2/}, \texttt{stma/anomaly\_det/event\_study/}, \texttt{stma/visualizations/impact\_analysis/}.
	extit{Time-series VAR / Granger track is documented separately in \texttt{VAR/} and is not included in this compendium.}}

\end{document}
